\section{Bài báo giải được gì, và để lại gì}
\label{sec:ch03-con-lai-gi}

\index{vanishing gradient!chưa giải được}%
Bài này lấy một quan sát đã có từ 1994 và biến nó thành hai điều kiện kiểm được
cùng hai biện pháp chạy được. Sau nó, câu hỏi \enquote{khi nào gradient tắt}
không còn là chuyện cảm tính. Đó là một bước tiến thật, và biện pháp thứ nhất
của bài vẫn còn nguyên trong PyTorch hôm nay dưới tên
\code{torch.nn.utils.clip\_grad\_norm\_}.

Nhưng hãy nhìn kỹ cái được sửa. Clipping chặn một bước đi quá dài. Nó làm huấn
luyện không hỏng. Nó không làm mạng nhớ được lâu hơn một bước nào: sau khi cắt,
gradient tại các bước xa vẫn nhỏ đúng như trước khi cắt, chỉ là toàn bộ vector
được co lại cùng một tỷ lệ. Clipping sửa được tính ổn định của huấn luyện, và
không sửa được trí nhớ.

Phía trí nhớ là việc của regularizer, và ở đó bài để lại một chỗ mắc kẹt mà
chính họ chỉ ra. Ép tín hiệu lỗi giữ nguyên norm khi đi ngược có nghĩa là đẩy
mô hình ra xa \tn{điểm hút}{attractor} mà nó đang tiến về, vì ở gần một điểm hút thì thông tin
buộc phải tắt dần. Đẩy ra xa điểm hút là đẩy lại gần ranh giới giữa các miền
hút. Mà vượt ranh giới giữa hai \tn{miền hút}{basin of attraction} chính là \tn{điều kiện đủ}{sufficient condition} để gradient nổ,
đúng như mục~\ref{sec:ch03-he-dong-luc} đã nói.

Hai biện pháp của bài vì thế không phải hai lựa chọn thay thế nhau mà là một
cặp buộc phải đi cùng: biện pháp xử lý phía tắt làm phía nổ dễ xảy ra hơn, nên
phải có biện pháp kia chặn lại. Bài gọi cấu hình đó là SGD-CR, và trên tập
test của cả ba bộ dữ liệu nhạc lẫn Penn Treebank, nó là cấu hình cho kết quả tốt
nhất~\autocite{pascanu2013difficulty}. Trên tập train thì không phải lúc nào cũng
vậy: ở ba trong bốn dòng, clipping một mình còn thấp hơn.

Đó là chỗ đứng mà bài để lại. Một \tn{mạng hồi quy}{recurrent neural network}
thường, để nhớ được xa, phải
được giữ ở một chế độ mà động lực học của nó không tự đi tới, bằng một số hạng
phạt kéo nó ra và một ngưỡng cắt chặn độ dài bước đi. Cân bằng ấy giữ được, và
bài chứng minh là giữ được. Chỉ là nó phải được giữ, liên tục, bằng thuật toán
huấn luyện.

Câu hỏi tự nhiên tiếp theo là câu hỏi mà chương sau trả lời: nếu đường đi của
gradient qua thời gian là thứ liên tục nhân với một cái gì đó nhỏ hơn 1, thì
thay vì phạt cho cái gì đó ấy lệch khỏi 1, có thể dựng một kiến trúc mà nó
\emph{bằng} 1 theo cấu tạo hay không.

Câu trả lời đã có từ 1997, mười sáu năm trước bài này. Nó không sửa thuật
toán huấn luyện. Nó sửa mạng.
